\chapter{Modèles de risque}\label{ch:rs:risk-models}

Un portefeuille neutre au marché est prévu à 3.3~\% de volatilité annuelle et réalise 9.2~\%. Son modèle de risque connaissait le marché, dix industries, le bêta, la taille et la valeur,
et le portefeuille était neutre à l'égard de tous ces facteurs~; il ne connaissait pas le momentum, et le score propriétaire du portefeuille, sans que personne ne l'ait vérifié, était pour moitié du momentum. Sur
huit ans, ses rendements standardisés ont un écart-type de 2.86 au lieu de un. Le même portefeuille, sous un modèle qui inclut le momentum, est prévu à
9.23~\% et réalise 9.20~\%. Un modèle de risque est un outil de recherche avant d'être un outil de reporting~: il dit quels risques un portefeuille prend, et il ne peut
seulement le dire que pour les risques qu'il connaît. Ce chapitre construit \texttt{firm.riskmodel}, un modèle fondamental d'actions dans le style que décrivent les éditeurs commerciaux
publiquement, l'ajuste à \texttt{firm.synthmkt} ainsi qu'un modèle statistique, et les évalue comme il faut, à l'aide de \omterm{def:rs:risk-models:bias}{statistiques de biais}.

\section{Trois familles de modèles factoriels}

Un modèle factoriel (livre~4, chapitre~22) écrit le rendement de chaque action comme les expositions multipliées par les \omterm{def:rs:risk-models:families}{rendements des facteurs}, plus un \omterm{def:rs:risk-models:families}{rendement spécifique}. Les modèles diffèrent par ce qu'ils observent.

\begin{definition}[Modèles à facteurs fondamentaux, statistiques et hybrides~; exposition aux facteurs, rendement du facteur, rendement spécifique]\label{def:rs:risk-models:families}
Dans un \emph{modèle à facteurs fondamentaux}\index{modèle à facteurs fondamentaux} les \emph{expositions aux facteurs}\index{expositions aux facteurs} sont des caractéristiques observées de chaque
action (appartenance à une industrie, scores de style standardisés) et les \emph{rendements des facteurs}\index{rendements des facteurs} sont estimés à chaque période par régression
en coupe transversale. Dans un \emph{modèle à facteurs statistiques}\index{modèle à facteurs statistiques}, les deux sont estimés à partir des seuls rendements, par composantes principales ou
analyse factorielle. Un \emph{modèle à facteurs hybride}\index{modèle à facteurs hybride} ajoute des facteurs statistiques à un modèle fondamental pour capturer ce que ses caractéristiques
manquent. Le \emph{rendement spécifique}\index{rendement spécifique} est ce que les facteurs laissent~: le résidu de la régression. (Une troisième famille, celle des modèles macroéconomiques, observe
les rendements des facteurs, comme l'inflation ou le spread de terme, et estime les expositions par régression en série temporelle.)
\end{definition}

\begin{definition}[Facteur de style, facteur sectoriel]\label{def:rs:risk-models:style}
Un \emph{facteur sectoriel}\index{facteur sectoriel} a pour exposition l'appartenance d'une action à une industrie~; un \emph{facteur de style}\index{facteur de style} a pour exposition une
caractéristique standardisée (taille, valeur, momentum, un bêta estimé), pondérée par la capitalisation pour avoir une moyenne nulle et ramenée à un écart-type en coupe transversale égal à un.
\end{definition}

Rosenberg a montré en 1974 que la covariance des rendements des actions comporte des composantes au-delà du marché et que leurs sensibilités suivent des caractéristiques observables~:
à l'origine du modèle fondamental. Connor a comparé les trois familles sur des rendements américains~; Chen, Roll et Ross ont évalué des facteurs macroéconomiques. Le modèle du chapitre
est fondamental, la famille que documentent les éditeurs commerciaux~: pour leur modèle américain, MSCI décrit un facteur pays qui rend les \omterm{def:rs:risk-models:style}{facteurs sectoriels} neutres en dollars
(chaque industrie mesurée nette du marché), un \omterm{def:rs:risk-models:specific}{risque spécifique} estimé à partir de séries temporelles quotidiennes et ramené vers la moyenne au sein de groupes de taille, et un \omterm{def:rs:risk-models:bias}{ajustement au régime de volatilité}.

\section{Estimation des rendements des facteurs}

Chaque jour $t$, le modèle régresse la coupe transversale des rendements sur les expositions connues à la clôture précédente~:
\[
\min_{f}\ \sum_i w_i\,(r_{it} - x_{i,t-1}^\top f)^2 \quad\text{sous la contrainte}\quad \sum_k c_k f_{k} = 0 \text{ sur les industries},
\]
avec des poids $w_i$ égaux à la racine carrée de la capitalisation (les grandes actions ont une variance spécifique plus petite) et $c_k$ la part de capitalisation de l'industrie $k$. Le
facteur pays a une exposition égale à un pour chaque action, si bien que les seules indicatrices sectorielles seraient colinéaires avec lui~; la contrainte fait du facteur pays le
marché pondéré par la capitalisation, et chaque \omterm{def:rs:risk-models:style}{facteur sectoriel} devient son rendement net du marché. \texttt{firm.riskmodel.cs\_regression} résout le problème contraint par
son système de Karush--Kuhn--Tucker (\cref{lst:rs:risk-models:csreg}).

Sur les dix années de \texttt{firm.synthmkt}, soit 1\,000 valeurs cotées par jour, le modèle compte quinze facteurs~: le facteur pays, dix \omterm{def:rs:risk-models:style}{facteurs sectoriels} et quatre \omterm{def:rs:risk-models:style}{facteurs de style}, chacun connu à
la clôture précédente~: le bêta, estimé sur l'année écoulée et réestimé chaque mois~; la taille, caractéristique latente de la simulation, fixée à la cotation~; la valeur,
le logarithme du rapport entre valeur comptable et cours, à partir de la dernière valeur comptable déposée (\texttt{firm.synthmkt.point\_in\_time\_styles})~; le momentum, rendement sur 12--1 mois, remis à zéro chaque
mois. La volatilité annuelle du facteur pays est de 17.8~\%, celle des industries entre 8.0~\% et 8.7~\%, celle du momentum de 7.2~\%. Le $R^2$ moyen quotidien en coupe transversale
est de 0.19~: l'essentiel de la variance quotidienne d'une action est spécifique, ce qui explique qu'un portefeuille de cent valeurs par côté conserve un \omterm{def:rs:risk-models:specific}{risque spécifique} qui mérite d'être modélisé.

Le respect de l'information disponible à chaque date compte ici comme partout. La simulation tient à jour, au fil de l'eau, les expositions de style de chaque action~; une première version du modèle de ce chapitre les
lisait à la fin de l'échantillon et les utilisait pour les dix années. Il s'ajustait sans sourciller, et son facteur momentum avait une volatilité de 3.3~\% au lieu de 7.2~\%~:
des expositions issues du futur décrivent un facteur différent, plus lisse.

\section{Covariance des facteurs}

La prévision de covariance des facteurs est une estimation à pondération exponentielle des \omterm{def:rs:risk-models:families}{rendements des facteurs}, avec une demi-vie de 90 jours (livre~4, chapitre~18, pour la
volatilité à pondération exponentielle). Pour les horizons de plus d'un jour, les termes de Newey--West ajoutent les autocovariances. Deux raffinements répondent aux deux
défaillances connues de l'estimateur. La volatilité n'est pas stationnaire~: une estimation rétrospective réagit avec retard quand la volatilité monte. Et l'erreur d'échantillonnage rend la covariance estimée trop
confiante sur certaines directions, que l'optimiseur exploite ensuite (Menchero, Wang et Orr~: le risque des portefeuilles optimisés est sous-estimé~; chapitre~25).

\begin{definition}[Statistique de biais, ajustement au régime de volatilité]\label{def:rs:risk-models:bias}
La \emph{statistique de biais}\index{statistique de biais} des prévisions de risque d'un portefeuille sur $T$ périodes est l'écart-type de ses rendements standardisés
$r_t/\hat\sigma_{t-1}$~; elle vaut un pour des prévisions correctes, et se situe dans la bande $1 \pm \sqrt{2/T}$ avec une probabilité d'environ 95~\% pour des rendements normaux. Le biais transversal du
jour $t$ est $B_t^2 = \frac1K\sum_k f_{kt}^2/\hat\sigma_{k,t-1}^2$ sur les facteurs. Un \emph{ajustement au régime de volatilité}\index{ajustement au régime de volatilité}
multiplie la prévision de covariance par une moyenne à pondération exponentielle $\lambda_t^2$ des $B_t^2$ récents, en relevant ensemble toutes les volatilités des facteurs lorsque les \omterm{def:rs:risk-models:families}{rendements des facteurs} du jour
sont grands par rapport à leurs prévisions.
\end{definition}

L'ajustement du chapitre utilise une demi-vie de 42 jours (\cref{lst:rs:risk-models:vra}). La volatilité du marché simulé suit un processus GJR-GARCH, si bien que les régimes
vont et viennent. Sur le portefeuille de marché pondéré par la capitalisation, la \omterm{def:rs:risk-models:bias}{statistique de biais} globale vaut 1.00 sans l'ajustement et 0.99 avec celui-ci~; la différence tient aux régimes.
Les \omterm{def:rs:risk-models:bias}{statistiques de biais} glissantes sur un an vont de 0.50 à 1.55 sans l'ajustement et de 0.56 à 1.32 avec lui, et 66~\% des fenêtres annuelles tombent hors de leur bande
($\pm 0.089$) sans l'ajustement, contre 44~\% avec lui (\cref{fig:rs:risk-models:rolling}). Une année reste longue pour un marché GARCH~; un desk qui gère son risque
au jour le jour veut l'ajustement et une demi-vie courte.

\section{Risque spécifique}

\begin{definition}[Risque spécifique]\label{def:rs:risk-models:specific}
Le \emph{risque spécifique}\index{risque spécifique} d'une action est l'écart-type prévu de son \omterm{def:rs:risk-models:families}{rendement spécifique}.
\end{definition}

La variance spécifique de chaque action est prévue par une moyenne à pondération exponentielle de ses résidus au carré (demi-vie de 90 jours). De telles prévisions sont bruitées, et
le bruit trie~: les actions dont la volatilité prévue est la plus forte sont en moyenne celles dont les résidus récents ont été par chance importants. Standardiser le résidu du lendemain de chaque action
par sa prévision et en regroupant par décile de la prévision, on le voit~: la \omterm{def:rs:risk-models:bias}{statistique de biais} vaut 1.13 dans le décile le plus bas et 0.93 dans le plus haut. Ramener chaque
prévision de volatilité de 10~\% du chemin vers la moyenne de son décile de taille, comme le décrit MSCI pour son modèle, aplanit la configuration pour la ramener entre 0.97 et 1.04
(\cref{fig:rs:risk-models:deciles}). Le niveau restant, légèrement supérieur à un, provient des chocs de Student-$t$ et des sauts aux dates de publication des résultats, présents dans la simulation.

\begin{omfigure}
\begin{tikzpicture}
\begin{axis}[width=10cm,height=5.2cm,ybar,bar width=7pt,xlabel={décile de la volatilité spécifique prévue},ylabel={statistique de biais},tick label style={font=\small},
  label style={font=\small},xmin=0.4,xmax=10.6,ymin=0.85,ymax=1.15,xtick={1,2,3,4,5,6,7,8,9,10},grid=major,grid style={gray!25},legend columns=2,
  legend style={font=\footnotesize,at={(0.5,-0.3)},anchor=north,draw=none,fill=none,/tikz/every even column/.append style={column sep=8pt}},
  table/col sep=comma]
\addplot[fill=gray!45,draw=gray] table[x=decile,y=raw]{figdata/research/24-risk-models/deciles.csv};
\addplot[fill=omDef!55,draw=omDef] table[x=decile,y=shrunk]{figdata/research/24-risk-models/deciles.csv};
\draw[black!60,densely dotted] (axis cs:0.4,1) -- (axis cs:10.6,1);
\legend{prévision en série temporelle,ramenée de 10~\% vers son décile de taille}
\end{axis}
\end{tikzpicture}

\omcaption{\omterm{def:rs:risk-models:bias}{Statistiques de biais} des \omterm{def:rs:risk-models:families}{rendements spécifiques} du lendemain standardisés par leurs prévisions, par décile de la prévision, années 3 à 10. Données~:
\texttt{rs\_riskmodel.specific\_deciles}.}\label{fig:rs:risk-models:deciles}
\end{omfigure}

\section{Évaluer un modèle de risque}

La variance prévue d'un portefeuille est $x^\top F x + \sum_i w_i^2 s_i^2$, avec $x = X^\top w$ ses \omterm{def:rs:risk-models:families}{expositions aux facteurs}, $F$ la covariance des facteurs et $s_i^2$ les
variances spécifiques~; la contribution $x_k (Fx)_k$ de chaque facteur se somme pour donner la partie factorielle (\cref{lst:rs:risk-models:risk}). Un modèle se juge sur la façon dont ses prévisions du risque de
nombreux portefeuilles se comparent à ce qui s'est passé, et non par la qualité de l'ajustement de sa régression.

Les portefeuilles de test du chapitre sont détenus de clôture à clôture sur le jour suivant et reconstruits chaque mois, et sont évalués sur les années 3 à 10 (2\,016 jours~; la bande de biais
vaut $1 \pm 0.031$). Cinquante portefeuilles aléatoires neutres au marché, cent valeurs par côté, ont des \omterm{def:rs:risk-models:bias}{statistiques de biais} comprises entre 0.96 et 1.04, avec une moyenne de 1.00~; 92~\% tombent dans la
bande (98~\% sans l'ajustement au régime, qui ne fait rien pour les portefeuilles dont le risque est surtout spécifique, et nuit un peu). Le portefeuille du début du chapitre est construit
à partir d'un score égal à 0.5 fois le score de momentum plus $\sqrt{0.75}$ fois du bruit, à l'achat les cent premières et à la vente les cent dernières, puis neutralisé vis-à-vis
de chaque facteur du modèle, sans le momentum. Le jour 1\,500, le modèle complet le prévoit à 6.88~\% par an, dont 76~\% de la variance provient de son exposition au momentum, égale à
1.03~; sur les huit années, son exposition a une racine carrée de la moyenne des carrés égale à 1.23.

\begin{center}
\small
\begin{tabular}{@{}lccc@{}}
\toprule
le portefeuille exposé, années 3 à 10 & prévision & réalisé & \omterm{def:rs:risk-models:bias}{statistique de biais}\\
\midrule
\omterm{def:rs:risk-models:families}{modèle à facteurs fondamentaux}, tous les facteurs & 9.23~\% & 9.20~\% & 1.04\\
\omterm{def:rs:risk-models:families}{modèle à facteurs fondamentaux} sans momentum & 3.26~\% & 9.20~\% & 2.86\\
\omterm{def:rs:risk-models:families}{modèle à facteurs statistiques}, 15 composantes principales & 5.59~\% & 9.20~\% & 1.68\\
\bottomrule
\end{tabular}
\end{center}

Le modèle sans momentum ne voit que le \omterm{def:rs:risk-models:specific}{risque spécifique} du portefeuille~; la volatilité de 7.2~\% du facteur momentum, multipliée par l'exposition du portefeuille, représente
les deux tiers manquants. Le modèle statistique, quinze composantes principales réestimées chaque mois sur l'année glissante, en retrouve une partie sans qu'on le lui dise (5.59~\%, une \omterm{def:rs:risk-models:bias}{statistique de biais}
de 1.68)~: les rendements du facteur momentum forment une direction de variation commune, que les composantes principales recherchent, mais ses chargements sont réinitialisés chaque mois
tandis que les composantes estimées sur l'année glissante supposent des chargements qui restent fixes, de sorte que la direction décrite par les données de l'année est une moyenne de douze directions différentes. Un
modèle statistique vaut encore la peine d'être utilisé à côté du modèle fondamental comme signal d'alarme~: ici il prévoit 5.59~\% là où le modèle fondamental annonce 3.26~\%, et quand le
risque du portefeuille diverge selon les deux modèles, le modèle fondamental omet quelque chose que détient le portefeuille. C'est aussi le cas pour les \omterm{def:rs:risk-models:families}{modèles à facteurs hybrides}.

\begin{omfigure}
\begin{tikzpicture}
\begin{axis}[name=a,width=6.8cm,height=5.4cm,xlabel={\small année},ylabel={\small statistique de biais sur un an},tick label style={font=\footnotesize},
  xmin=3,xmax=10,ymin=0.4,ymax=3.6,grid=major,grid style={gray!25},title={\small le portefeuille exposé},legend columns=1,
  legend style={font=\scriptsize,at={(0.5,-0.32)},anchor=north,draw=none,fill=none},table/col sep=comma]
\addplot[omDef,very thick] table[x=year,y=full]{figdata/research/24-risk-models/rolling.csv};
\addplot[omThm,thick] table[x=year,y=nomom]{figdata/research/24-risk-models/rolling.csv};
\addplot[black!60,thick,dashed] table[x=year,y=pca]{figdata/research/24-risk-models/rolling.csv};
\draw[black!50,densely dotted] (axis cs:3,1) -- (axis cs:10,1);
\legend{tous les facteurs,sans momentum,statistique}
\end{axis}
\begin{axis}[at={(a.east)},anchor=west,xshift=1.2cm,width=6.8cm,height=5.4cm,xlabel={\small année},tick label style={font=\footnotesize},
  xmin=3,xmax=10,ymin=0.4,ymax=1.8,grid=major,grid style={gray!25},title={\small le portefeuille de marché},legend columns=1,
  legend style={font=\scriptsize,at={(0.5,-0.32)},anchor=north,draw=none,fill=none},table/col sep=comma]
\addplot[black!60,thick] table[x=year,y=mkt_raw]{figdata/research/24-risk-models/rolling.csv};
\addplot[omDef,very thick] table[x=year,y=mkt_vra]{figdata/research/24-risk-models/rolling.csv};
\draw[black!50,densely dotted] (axis cs:3,1) -- (axis cs:10,1);
\legend{sans ajustement,ajustement au régime}
\end{axis}
\end{tikzpicture}

\omcaption{\omterm{def:rs:risk-models:bias}{Statistiques de biais} glissantes sur un an. À gauche~: le portefeuille exposé au momentum selon les trois modèles. À droite~: le marché pondéré par la capitalisation selon le modèle complet,
avec et sans l'\omterm{def:rs:risk-models:bias}{ajustement au régime de volatilité}. Données~: \texttt{rs\_riskmodel}.}\label{fig:rs:risk-models:rolling}
\end{omfigure}

Les \omterm{def:rs:risk-models:bias}{statistiques de biais} ont leur propre incertitude. Sur soixante mois, la bande vaut $1 \pm 0.18$~; avec des queues épaisses, moins de 95~\% des prévisions justes tombent à l'intérieur
(MSCI rapporte 86~\% pour une kurtosis de 5 sur 120 périodes). La \omterm{def:rs:risk-models:bias}{statistique de biais} d'un seul portefeuille est un indice faible~; de nombreux portefeuilles, fenêtres glissantes et les
portefeuilles que la firme négocie réellement sont ce sur quoi un modèle est jugé.

\section{Tutoriel~: prévu trois, réalisé neuf}\label{tut:rs:risk-models:tut}

\begin{tutorial}
\textbf{Objectif.} Ajuster un \omterm{def:rs:risk-models:families}{modèle à facteurs fondamentaux} avec et sans momentum, ainsi qu'un modèle statistique, sur \texttt{firm.synthmkt}, prévoir le risque de portefeuilles \omterm{def:rs:overfitting:overfitting}{hors échantillon},
et les évaluer par \omterm{def:rs:risk-models:bias}{statistiques de biais}. \textbf{État final~:} le tableau et les \cref{fig:rs:risk-models:rolling,fig:rs:risk-models:deciles}.
\begin{tutsteps}
\item \textbf{La régression quotidienne} avec la contrainte sectorielle.
\omcode{code/firm/riskmodel/firm_riskmodel.py}{30}{47}{Régression transversale pondérée sous contrainte.}\label{lst:rs:risk-models:csreg}
\item \textbf{L'\omterm{def:rs:risk-models:bias}{ajustement au régime de volatilité}.}
\omcode{code/firm/riskmodel/firm_riskmodel.py}{87}{98}{Le multiplicateur issu de la statistique de biais transversale.}\label{lst:rs:risk-models:vra}
\item \textbf{Risque du portefeuille} et sa décomposition.
\omcode{code/firm/riskmodel/firm_riskmodel.py}{128}{134}{Variance factorielle et spécifique, et contribution de chaque facteur.}\label{lst:rs:risk-models:risk}
\item \textbf{Exécuter} \texttt{rs\_riskmodel.fundamental()}, \texttt{fundamental} sans le momentum, \texttt{statistical()}, \texttt{summary()}, \texttt{random\_bias()},
\texttt{market\_bias()}, \texttt{specific\_deciles()} et \texttt{fig\_riskmodel.py}.
\end{tutsteps}
\textbf{Pour aller plus loin.} Retirer le facteur valeur au lieu du momentum et construire un portefeuille exposé à ce facteur~; ajouter trois composantes principales des résidus du modèle fondamental
comme facteurs statistiques (un \omterm{def:rs:risk-models:families}{modèle à facteurs hybride}) et voir si le facteur manquant réapparaît.
\end{tutorial}

\section{Construction~: un modèle de risque actions}\label{bld:rs:risk-models:riskmodel}

\begin{build}
\bfield{Rôle} Le modèle de risque actions de la firme~: expositions, \omterm{def:rs:risk-models:families}{rendements des facteurs}, covariance et \omterm{def:rs:risk-models:specific}{risque spécifique} chaque jour, le risque de tout portefeuille et sa décomposition,
et les éléments (\omterm{def:rs:risk-models:bias}{statistiques de biais}) montrant que les prévisions sont justes.
\bfield{Interface} \texttt{cs\_regression(r, X, w, C)}, \texttt{factor\_returns(R, exposures, W, C)}, \texttt{ewma\_cov(F, half\_life, nw\_lags)},
\texttt{vra(F, covs, half\_life)}, \texttt{specific\_var(E, half\_life, groups, shrink)}, \texttt{portfolio\_risk(w, X, Fcov, spec)}, \texttt{pca\_model(R, k)},
\texttt{bias\_stat(r, sigma)}, \texttt{bias\_band(T)}.
\bfield{Règles} Expositions connues à la clôture précédente~; prévisions établies à la clôture pour la période suivante~; tout changement de modèle est jugé sur les \omterm{def:rs:risk-models:bias}{statistiques de biais} de
portefeuilles aléatoires, des portefeuilles factoriels et des portefeuilles que détient la firme~; un modèle statistique tourne à côté du modèle fondamental.
\bfield{Tests d'acceptation} \texttt{code/firm/riskmodel/tests/}~: une régression contrainte qui retrouve les rendements de facteurs injectés et respecte sa contrainte~; la
covariance et le multiplicateur de régime détectant un décalage de volatilité injecté~; les variances spécifiques retrouvant un niveau injecté, le rétrécissement réduisant leur dispersion~;
contributions dont la somme égale la variance factorielle~; composantes principales orthonormées~; \omterm{def:rs:risk-models:bias}{statistiques de biais} de prévisions connues et de la bande.
\bfield{Pour aller plus loin} L'ajustement eigenfactor du risque des portefeuilles optimisés~; facteurs statistiques hybrides~; horizons multi-jours avec termes de Newey--West.
\end{build}

\begin{omsources}
\item B.~Rosenberg, ``Extra-market components of covariance in security returns'', \emph{Journal of Financial and Quantitative Analysis} 9(2), 1974.
\item G.~Connor, ``The three types of factor models: a comparison of their explanatory power'', \emph{Financial Analysts Journal} 51(3), 1995.
\item N.-F.~Chen, R.~Roll and S.~A.~Ross, ``Economic forces and the stock market'', \emph{Journal of Business} 59(3), 1986.
\item Y.~Liu, J.~Menchero, D.~J.~Orr and J.~Wang, \emph{The Barra US Equity Model (USE4): Empirical Notes}, MSCI, 2011.
\item J.~Menchero, J.~Wang and D.~J.~Orr, ``Improving risk forecasts for optimized portfolios'', \emph{Financial Analysts Journal} 68(3), 2012.
\end{omsources}

\section{Exercices}

\begin{exercise}[$\star$]\label{exo:rs:risk-models:1}
La volatilité prévue du portefeuille, due au seul \omterm{def:rs:risk-models:specific}{risque spécifique}, est de 3.26~\%. Son exposition à un facteur omis, de volatilité 7.2~\%, a une racine carrée de la moyenne des carrés égale à 1.23. Déterminer
la volatilité réalisée qui en résulte et la \omterm{def:rs:risk-models:bias}{statistique de biais} correspondante.
\end{exercise}

\begin{exercise}[$\star$]\label{exo:rs:risk-models:2}
Donner la bande à 95~\% de la \omterm{def:rs:risk-models:bias}{statistique de biais} pour 60 observations mensuelles et pour 2\,016 observations quotidiennes.
\end{exercise}

\begin{exercise}[$\star$]\label{exo:rs:risk-models:3}
Expliquer pourquoi un modèle comportant un facteur pays et un ensemble complet de variables indicatrices sectorielles nécessite une contrainte, et ce que la contrainte pondérée par la capitalisation fait signifier aux facteurs.
\end{exercise}

\begin{exercise}[$\star\star$]\label{exo:rs:risk-models:4}
Expliquer pourquoi les prévisions de séries temporelles de la volatilité spécifique surestiment les valeurs à forte volatilité et sous-estiment les valeurs à faible volatilité, même
lorsque la volatilité de chaque valeur est constante.
\end{exercise}

\begin{exercise}[$\star\star$]\label{exo:rs:risk-models:5}
L'ajustement au régime a amélioré les \omterm{def:rs:risk-models:bias}{statistiques de biais} glissantes du portefeuille de marché et a légèrement dégradé celles des portefeuilles aléatoires. Pourquoi~?
\end{exercise}

\begin{exercise}[$\star\star$]\label{exo:rs:risk-models:6}
Montrer que les contributions des facteurs $x_k(Fx)_k$ somment à la variance du facteur. Une contribution peut-elle être négative~?
\end{exercise}

\begin{exercise}[$\star\star\star$]\label{exo:rs:risk-models:7}
\emph{Programmation.} Construire le portefeuille exposé à partir d'un score de charge 0.3 au lieu de 0.5 sur le momentum, et rapporter les \omterm{def:rs:risk-models:bias}{statistiques de biais} des trois modèles.
\end{exercise}

\begin{exercise}[$\star\star\star$]\label{exo:rs:risk-models:8}
\emph{Trouver la faille.} «~Notre modèle de risque n'explique que 19~\% de la variance quotidienne des rendements, il est donc inutile pour un portefeuille de deux cents noms.~»
\end{exercise}

\section{Problème~: Prédit trois, réalisé neuf}

\begin{problem}[{Problème du week-end --- un modèle qui ne savait pas}]\label{pb:rs:risk-models:1}
Les modèles du chapitre sur \texttt{firm.synthmkt}.

\medskip\noindent\textbf{Partie I --- Le modèle.}
\begin{enumerate}
\item Lister les quinze facteurs et la manière dont chaque exposition est obtenue.
\item Pourquoi des pondérations en racine carrée de la capitalisation~?
\item Quelles sont les volatilités des facteurs pays, sectoriels et du momentum~?
\item Quelle est la moyenne transversale du $R^2$, et que signifie-t-elle pour le risque d'un portefeuille~?
\end{enumerate}

\medskip\noindent\textbf{Partie II --- Covariance et \omterm{def:rs:risk-models:specific}{risque spécifique}.}
\begin{enumerate}[resume]
\item Comment fonctionne l'\omterm{def:rs:risk-models:bias}{ajustement au régime de volatilité}, et que fait-il aux \omterm{def:rs:risk-models:bias}{statistiques de biais} glissantes du portefeuille de marché~?
\item Pourquoi n'aide-t-il pas les portefeuilles aléatoires~?
\item Quel profil les prévisions spécifiques brutes montrent-elles par décile, et que fait la contraction~?
\item Pourquoi subsiste-t-il un niveau résiduel légèrement supérieur à un~?
\end{enumerate}

\medskip\noindent\textbf{Partie III --- Le portefeuille exposé.}
\begin{enumerate}[resume]
\item Comment le portefeuille exposé est-il construit~?
\item Quelles sont ses prévisions, sa volatilité réalisée et ses \omterm{def:rs:risk-models:bias}{statistiques de biais} sous les trois modèles~?
\item Quelle part de sa variance prévue provient du momentum, et à quelle exposition~?
\item Pourquoi le modèle statistique retrouve-t-il la plus grande part du risque manquant, mais pas la totalité~?
\end{enumerate}

\medskip\noindent\textbf{Partie IV --- Le verdict.}
\begin{enumerate}[resume]
\item Énoncer le \emph{résultat nommé}~: la \omterm{def:rs:risk-models:bias}{statistique de biais} du modèle qui ignore le facteur introduit, et le rapport entre volatilité réalisée et volatilité prévue du portefeuille exposé.
book.
\item Comment le desk l'aurait-il découvert avant le drawdown~?
\item Quelles preuves exigeriez-vous avant d'adopter un nouveau modèle de risque~?
\item Comment les queues épaisses modifient-elles la lecture d'une \omterm{def:rs:risk-models:bias}{statistique de biais}~?
\item Qu'est-ce qu'un \omterm{def:rs:risk-models:families}{modèle à facteurs hybride}, et quand vaut-il la peine d'être construit~?
\item Pourquoi les \omterm{def:rs:risk-models:bias}{statistiques de biais} des portefeuilles aléatoires constituent-elles un test nécessaire mais insuffisant~?
\item Quel portefeuille un gestionnaire des risques doit-il tester en premier~?
\item En une phrase~: que sait un modèle de risque~?
\end{enumerate}
\end{problem}

\section{Questions d'entretien}

\begin{interviewq}[$\star$ \iqroles{risque, chercheur}]\label{iq:rs:risk-models:1}
Quelle est la différence entre un \omterm{def:rs:risk-models:families}{modèle à facteurs fondamentaux} et un \omterm{def:rs:risk-models:families}{modèle à facteurs statistiques}~? Quand préféreriez-vous chacun~?
\end{interviewq}

\begin{interviewq}[$\star\star$ \iqroles{chercheur}]\label{iq:rs:risk-models:2}
Comment les \omterm{def:rs:risk-models:families}{rendements des facteurs} d'un modèle fondamental sont-ils estimés chaque jour, et pourquoi la contrainte sur les secteurs~?
\end{interviewq}

\begin{interviewq}[$\star\star$ \iqroles{risque}]\label{iq:rs:risk-models:3}
Un portefeuille neutre au marché réalise deux fois sa volatilité prévue. Comment découvririez-vous la cause~?
\end{interviewq}

\begin{interviewq}[$\star\star$ \iqroles{risque, chercheur}]\label{iq:rs:risk-models:4}
Comment évaluez-vous un modèle de risque~? Qu'est-ce qu'une \omterm{def:rs:risk-models:bias}{statistique de biais}, et quels en sont les pièges~?
\end{interviewq}

\begin{interviewq}[$\star\star$ \iqroles{chercheur}]\label{iq:rs:risk-models:5}
Pourquoi les portefeuilles optimisés ont-ils tendance à voir leur risque sous-estimé~?
\end{interviewq}

\begin{interviewq}[$\star\star\star$ \iqroles{chercheur, développeur}]\label{iq:rs:risk-models:6}
Concevez la production quotidienne d'un modèle de risque actions~: entrées, ordre de calcul, contrôles et données stockées.
\end{interviewq}
